\begin{DatosTabla}{L{0.24\linewidth} C{0.14\linewidth} L{0.44\linewidth}}{\textbf{Medida} & \textbf{Valor} & \textbf{Interpretación}}[Analisis descriptivo de math\_score]

Moda & 85 & Para la variable \Variable{math\_score}, el valor que más se repite es 85. \\ \hline
Media & 85,09 & En promedio los alumnos obtuvieron 85,09 como puntaje de matemática. \\ \hline
Mediana & 82 & El 50\% de los estudiantes sacó 82 o menos. \\ \hline
Primer Cuartil & 80 & El 25\% de los estudiantes sacó 80 o menos. \\ \hline
Tercer Cuartil & 85 & El 75\% de los estudiantes sacó 85 o menos. \\ \hline
Varianza & 10,62 & Los puntajes presentan una dispersión baja, ya que la mayoría de los estudiantes obtuvo resultados cercanos al promedio de 85,09 puntos. \\ \hline
Desvío Estándar & 3,26 & Los puntajes suelen diferir de la media en aproximadamente 3,26 puntos. \\ \hline
Rango Intercuartílico & 5,00 & En el 50\% central hay una diferencia de 5 puntos. \\ \hline
Coeficiente de Variación & 3,83\% & Como es menor al 20\%, se puede decir que los datos son homogéneos y, por lo tanto, la media es representativa de este conjunto de datos. \\


\end{DatosTabla}

